% 第8章 生成式模型
\section{生成式模型：合成数据、场景生成与市场仿真}\label{sec:gen}

生成式模型在这一族里的角色和前几族不同：它不直接预测收益，也不直接做决策，而是\textbf{生产数据}。这个角色看似边缘，实则在交易链条上占着一个战略位置——因为前面几族都需要“样本”，而金融恰恰是最缺样本的领域。合成数据、场景生成与市场仿真，本质都是在试图人造样本，以缓解“历史只有一份”这条硬约束。

\subsection{问题设定与三种用途}

生成式模型（GAN、VAE、扩散模型）学习数据分布 $p_\text{data}$ 并采样出新样本。在交易场景里，它有三种用法：

\begin{enumerate}
  \item \textbf{数据增强}：为罕见事件（极端行情、流动性枯竭）或样本稀薄的资产/市场生成补充样本。
  \item \textbf{场景与压力测试}：生成“可能发生但历史上还没发生”的市场路径，用于评估策略与组合的尾部表现。
  \item \textbf{市场仿真}：构建一个能对交易行为作出反应的虚拟市场，为强化学习提供训练环境（第 \ref{sec:rl} 章反复强调“模拟器就是一切”，生成式模型正是造模拟器的主要工具）。
\end{enumerate}

\subsection{方法与判据：什么叫“像”}

金融时序生成有一把特殊的尺子，叫\textbf{风格化事实}（stylized facts）：收益的肥尾、波动率聚集、杠杆效应、绝对收益的长记忆、成交量的相关性等。一个好生成器必须在统计上复现这些事实，而不仅仅是“目测像”。

\textbf{GAN 一代}。Wiese 等（2020）\cite{wiese2020quantgan} 提出的 Quant GANs 用（时序）卷积结构生成金融路径，并系统地以风格化事实作为验收标准，是这一方向被引用最多的工作之一。Takahashi 等（2019）\cite{takahashi2019} 用 GAN 建模金融时序。Yoon、Jarrett 与 van der Schaar（2019）\cite{yoon2019timegan} 的 TimeGAN 提出了“监督式”生成，把时序的动态与静态特征分开建模，是通用时序生成的代表。这一代的共同局限是训练不稳定、模式坍缩（只会生成少数几种路径），对齐尾部与波动率聚集尤其困难。

\textbf{扩散模型一代}。2023 年后，扩散模型因其训练更稳定、覆盖分布更广，被引入金融时序与订单流生成。相对于 GAN，扩散模型在复现尾部与多模态分布上更有优势，但采样成本高。这一代目前在交易上的公开、可复现证据仍有限。

\textbf{与对冲/执行的结合}。Buehler 等（2019）\cite{buehler2019deephedging} 的 Deep Hedging 把生成模型与对冲策略联合训练：先生成市场路径，再在该分布上学一个对冲策略。它的价值在于把“不确定性建模”与“决策优化”接在一起，是生成式模型进入交易决策的干净范例。

\textbf{市场仿真与智能体}。另一条路是自底向上的仿真：用多智能体（agent-based）模型复现市场的涌现行为，ABIDES（Byrd 等 2020\cite{byrd2020abides}）是较完整的开源离散事件市场模拟器。2023 年后，“生成式智能体”（Park 等 2023\cite{park2023generative}）的思路被迁移到市场仿真，用 LLM 扮演各类投资者，试图生成更真实的市场微结构。这类工作前景诱人，但目前主要是“生成更逼真的行为”，把它作为交易策略的可信验证环境还为时过早。

\subsection{生成质量的评估清单}

“像不像”需要一把可核查的尺子。参照文献里较通用的做法，本文把金融序列生成的质量评估归成四层，越往下越接近交易实用：

\begin{enumerate}
  \item \textbf{边缘分布层}：收益的均值、方差、偏度与非正态（肥尾）；波动率（绝对收益）的自相关是否复现。这一层是必要条件，但很容易通过。
  \item \textbf{联合结构层}：跨资产的相关系数矩阵、尾部相依（极端行情下的相关性上升）、波动率聚集与杠杆效应。这一层是多数生成器的分水岭。
  \item \textbf{路径层}：长记忆性、跳跃与波动率聚集的时序形态、是否出现“平滑得不真实”的路径。GAN 的模式坍缩常在此暴露。
  \item \textbf{下游任务层}：用合成数据训练一个策略或风险模型，看它在\textbf{真实数据}上的表现与用真实数据训练相比如何。这是唯一有交易意义的验收，却最常被跳过。
\end{enumerate}

四层里，前三层可以自动化（一套统计检验即可），第四层需要真实数据的对照实验。一个常见误区是把前三层当成全部，于是得到一个“统计上很像、训练起来没用”的生成器。

\subsection{生成器当训练环境时的耦合风险}

生成式模型最诱人的用法，是“造一个环境给强化学习练”（第 \ref{sec:rl} 章）。这条路有一个结构性风险：\textbf{策略会和生成器一起过拟合到生成器自身的偏差上}——生成器造不出的场景，策略自然学不会；生成器造得过多的场景，策略会学得过重。等到接上真实市场，两者一起失效。

缓解办法有三条，都不是新技术，而是纪律：

\begin{enumerate}
  \item \textbf{生成器与策略交替更新，而不是一次冻结}。生成器一旦冻结，退化成了固定的历史回放；交替训练至少让环境随策略演化，逼近对抗关系。
  \item \textbf{保留一段真实数据做最终裁判}。用它训练的策略都在合成环境里调参，但最终验收必须在未参与训练的真实留出段上完成，且口径与 P0--P2 一致。
  \item \textbf{报告“分布外表现”}。明确给出策略在生成器覆盖范围之外（更极端波动、更低流动性）的表现；如果只在训练分布内有效，应直接判为“生成器偏差”，而不是“策略有效”。
\end{enumerate}

一句话总结这条纪律：\textbf{用合成环境训练可以提高样本效率，但不能用它来证明策略有效；证明的那一步，永远要在真实数据上、用第 16 章的口径完成。}

\subsection{落地要素}

\begin{itemize}
  \item \textbf{保真度评估要下沉到下游任务}。风格化事实通过了，不代表用它训练的策略就有效。必须做“下游任务保真”：用合成数据训练一个策略，看它在真实数据上的表现，而这一步恰恰最常被省略。
  \item \textbf{绝不能“用生成数据证明策略有效”}。如果策略在合成数据上表现好，只能说明策略适配了这个生成器的假设，不能说明它在真实市场有效。这是评估合成数据价值的红线。
  \item \textbf{尾部与相关性优先}。对交易而言，均值和方差容易对，难的是尾部与跨资产的依赖结构；而后者恰恰是风险管理的命脉。
\end{itemize}

\subsection{复现状态与失效边界}

Quant GANs\cite{wiese2020quantgan} 与 TimeGAN\cite{yoon2019timegan} 有公开实现，复现状态为\textbf{部分复现}；扩散模型在金融生成上的工作复现状态普遍为\textbf{待复核}（缺代码、缺标准数据集）。失效边界最关键的一条是\textbf{“只学历史分布”}：生成模型的训练集是历史，它能造出的“极端场景”，其极端程度不可能超过它见过的范围——用一个在平静期训练的生成器做压力测试，会系统性低估风险，因为它根本没学过平静期之外的世界。这构成了一个悖论式的局限：\textbf{你最需要合成数据的时刻（从未发生过的危机），恰恰是生成模型最不可信的时刻。}

本章要记住的一句话是：\textbf{生成式模型缓解了“历史只有一份”的样本焦虑，却无法制造真正的、未曾发生过的历史——它能把已知的分布采样得更好，但把未知当作已知来生成，正是最危险的那类错误。}
